\chapter{The Weight Map and Lagrangian Structure}
\label{sec:weighted}
% ===================================================================

\section{Weighted GSLTs}

\begin{definition}[Weight Map]
A \emph{weight map} for a GSLT $\GSLT$ is a family of functions
\[
  \fwd_r : \HML(\ctx) \to \CC, \quad r \in \rules
\]
where the domain of $\fwd_r$ consists of formulae that \emph{refine the
left-hand side of $r$}: formulae $\phi$ such that $u \sat \phi$ implies $u$
matches the pattern $l_r$.  The weight $\fwd_r(\phi) \in \CC$ is the
\emph{amplitude} associated with firing rule $r$ on a term satisfying $\phi$.
\end{definition}

\begin{definition}[Weighted GSLT]
A \emph{weighted GSLT} is a pair $(\GSLT, \fwd)$ where $\GSLT$ is a GSLT and
$\fwd = (\fwd_r)_{r \in \rules}$ is a weight map.
\end{definition}

\section{The Lagrangian}

The weight map equips every rewrite path with an amplitude.

\begin{definition}[Path Amplitude]
Let $\gamma = r_1 \cdot r_2 \cdots r_n$ be a rewrite path from $P$ to $Q$,
and let $\phi_i$ be the most refined formula in the domain of $\fwd_{r_i}$
satisfied by the term at step $i$.  The \emph{path amplitude} is:
\[
  W(\gamma) \;=\; \prod_{i=1}^{n} \fwd_{r_i}(\phi_i) \;\in\; \CC.
\]
\end{definition}

\begin{definition}[Action Functional]
The \emph{action} of a path $\gamma$ is:
\[
  S[\gamma] \;=\; \sum_{i=1}^{n} \log \fwd_{r_i}(\phi_i)
\]
when all weights are nonzero, so that $W(\gamma) = e^{S[\gamma]}$.
In the real-valued special case $\fwd_r : \HML(\ctx) \to \RR_{>0}$, this
is a standard additive action functional.
\end{definition}

\begin{remark}[The Lagrangian as a Step Cost Function]
The Lagrangian of classical mechanics is a function $L(q, \dot{q}, t)$
assigning a real number to each instantaneous state-and-velocity.  In the
present framework, $\fwd_r(\phi)$ plays this role: it assigns an amplitude
to an instantaneous term and its ``velocity'' (the rule $r$ it is about to
fire, classified by the formula $\phi$ it satisfies).  The path amplitude
$W(\gamma)$ is the multiplicative analogue of $e^{iS/\hbar}$ in the Feynman
path integral, and the principle of stationary action corresponds to the
saddle-point approximation of $\sum_\gamma W(\gamma)$.
\end{remark}

\begin{remark}[The Uniform Case]
When $\fwd_r(\phi) = 1$ for all $r, \phi$, every path has unit amplitude and
$d_{\min}(P, Q)$ counts the number of paths of minimum length.  The path
integral degenerates to counting minimal causal histories.
\end{remark}

% ===================================================================
